% Original vector diagram for source Fig.29.2, PDF129 / printed117;
% its explanatory paragraph is on PDF130 / printed118.
% The manuscript supplies figure label C29b and its caption.
% One six-valent insertion: one high-mode self-loop uses two half-edges;
% four low-mode external legs remain. The legend names the action coefficient,
% not the Euclidean vertex factor, which is -6! c_{6,1}(Lambda_0).
\begin{tikzpicture}[x=12mm,y=12mm,line width=.7pt,
  font=\fontsize{10.5}{12}\selectfont]
  \coordinate (v) at (0,0);
  \draw[dashed] (v)
    .. controls (-.75,.55) and (-1.05,1.3) .. (-.55,1.75)
    .. controls (-.20,2.10) and (.20,2.10) .. (.55,1.75)
    .. controls (1.05,1.3) and (.75,.55) .. (v);
  \draw (v) -- (-1.75,0);
  \draw (v) -- (1.75,0);
  \draw (v) -- (-1.20,-1.2);
  \draw (v) -- (1.20,-1.2);
  \fill (v) circle[radius=.65mm];
  \node[left] at (-1.75,0) {$p_1$};
  \node[right] at (1.75,0) {$p_2$};
  \node[below left] at (-1.20,-1.2) {$p_3$};
  \node[below right] at (1.20,-1.2) {$p_4$};
  \node at (0,1.28) {$\ell$};
  \draw (2.8,1.2) -- (3.5,1.2);
  \node[right] at (3.67,1.2) {低动量：$|p_i|<\Lambda$};
  \draw[dashed] (2.8,.61) -- (3.5,.61);
  \node[right] at (3.67,.61) {$\Lambda<|\ell|<\Lambda_0$};
  \fill (3.15,-.10) circle[radius=.65mm];
  \node[right] at (3.67,-.10) {$c_{6,1}(\Lambda_0)\varphi^6$};
\end{tikzpicture}
